\chapter{Recursividad y relaciones recurrentes}
\label{md:recurrencias}

\section{Definición recursiva y condiciones iniciales}

Una definición recursiva parte de elementos explícitos y construye
los restantes a partir de otros ya obtenidos.
Para una sucesión, una \textbf{relación de recurrencia} expresa un
término mediante términos anteriores; los \textbf{valores iniciales}
determinan cuál de las sucesiones que satisfacen la relación se elige.
La relación aislada generalmente no determina una única sucesión.

El factorial, por ejemplo, se fija mediante $0!=1$ y $n!=n(n-1)!$.
Los árboles con raíz también admiten construcción recursiva:
se parte de un único vértice y se une una nueva raíz a las raíces de
árboles disjuntos ya construidos.
La inducción del capítulo~\ref{md:induccion} permite comprobar
propiedades de los objetos definidos así.

\section{Cómo modelar un recuento recursivo}
\label{md:modelar-recurrencia}

Se divide el conjunto de objetos según una decisión que deje subobjetos
más pequeños. Los casos deben ser excluyentes y exhaustivos;
se aplican suma y producto para relacionar sus cantidades.
Después se calculan los valores iniciales para los tamaños
en los que la descomposición todavía no se aplica.

Para las desordenaciones de la sección~\ref{md:desordenaciones},
\[
 d(1)=0,\quad d(2)=1,\qquad
 d(n)=(n-1)\bigl(d(n-1)+d(n-2)\bigr)\quad(n\ge3).
\]
La idea es elegir $a=\sigma(1)$, con $n-1$ opciones.
Si $\sigma(a)=1$, se desordena el resto de $n-2$ elementos.
En caso contrario, se elimina $1$ de la permutación, enlazando
su predecesor directamente con $a$; se obtiene una desordenación
de $n-1$ elementos, y el proceso se puede invertir.

\section{Recurrencias lineales con coeficientes constantes}
\label{md:caracteristica}

Una relación de la forma
\[
 r(n)=a_1r(n-1)+\cdots+a_tr(n-t)+h(n),\qquad n\ge t+1,
\]
es \textbf{lineal con coeficientes constantes} si los $a_i$ son constantes.
Su orden es $t$ cuando $a_t\ne0$. Es \textbf{homogénea} si $h(n)=0$.
Los productos entre términos de la sucesión y las potencias de dichos
términos no corresponden a este modelo lineal.
Dados $r(1),\ldots,r(t)$, la relación determina los demás valores
sucesivamente.

Para la recurrencia homogénea, la \textbf{ecuación característica} es
\[
 x^t-a_1x^{t-1}-\cdots-a_{t-1}x-a_t=0.
\]
Para $b\ne0$, la sucesión $r(n)=b^n$ satisface la recurrencia
si y solo si $b$ es raíz de esa ecuación: se sustituye y se divide
por $b^{n-t}$.

Cuando hay $t$ raíces reales distintas y no nulas $b_1,\ldots,b_t$,
todas las soluciones tienen la forma
\[
 \boxed{r(n)=c_1b_1^n+\cdots+c_tb_t^n.}
\]
Las combinaciones lineales conservan la relación homogénea.
Las constantes se obtienen imponiendo los $t$ valores iniciales:
\[
 \sum_{j=1}^t c_jb_j^i=r(i),\qquad i=1,\ldots,t.
\]
Este sistema tiene solución única: su determinante es un producto
de raíces no nulas por un determinante de Vandermonde, no nulo al
ser distintas. Al coincidir los primeros términos, la recurrencia
asegura por inducción que coincide toda la sucesión.

El procedimiento es: identificar el orden y la homogeneidad,
formar la ecuación, comprobar las condiciones sobre sus raíces,
escribir la combinación y resolver las constantes con los datos iniciales.
La fórmula anterior tiene ese alcance; no se aplica directamente
a raíces múltiples ni a una recurrencia no homogénea.

\begin{ejemplo}
Para $r(n)=3r(n-1)-2r(n-2)$, $r(1)=2$, $r(2)=5$,
la ecuación es $(x-1)(x-2)=0$.
De $r(n)=c_1+c_2 2^n$ y las dos condiciones iniciales se obtiene
$c_1=-1$, $c_2=3/2$, luego $r(n)=3\cdot2^{n-1}-1$.
La recurrencia y los datos iniciales quedan ambos satisfechos.
\end{ejemplo}

\section{Sucesión de Fibonacci}
\label{md:fibonacci}

Con la indexación
\[
 \operatorname{fib}(1)=1,\qquad \operatorname{fib}(2)=2,\qquad
 \operatorname{fib}(n)=\operatorname{fib}(n-1)+\operatorname{fib}(n-2)
 \quad(n\ge3),
\]
los primeros términos son $1,2,3,5,8,\ldots$.
Una \textbf{recurrencia de tipo Fibonacci} tiene esa misma relación,
pero puede tener otros valores iniciales.

La ecuación característica es $x^2-x-1=0$, con raíces
\[
 \alpha=\frac{1+\sqrt5}{2},\qquad
 \beta=\frac{1-\sqrt5}{2}.
\]
Al imponer los valores iniciales se obtiene la \textbf{fórmula de Binet}:
\[
 \boxed{\operatorname{fib}(n)=
  \frac{\alpha^{n+1}-\beta^{n+1}}{\sqrt5}\quad(n\ge1).}
\]
El exponente $n+1$ corresponde a los datos iniciales $1,2$.
Como $|\beta|<1$, el término $\beta^{n+1}/\sqrt5$ tiende a cero,
y $\alpha^{n+1}/\sqrt5$ aproxima la sucesión para índices grandes.
